\documentclass{article}
\usepackage{amsmath}
\usepackage{graphicx}
\graphicspath{ {images/} }
\usepackage[spanish]{babel}
\usepackage{bbm}
\usepackage{amsthm}
\usepackage{authblk}
\usepackage{hyperref}
\usepackage[utf8]{inputenc}
\usepackage{mathrsfs}
\usepackage{amsfonts}
\usepackage{amssymb}
\usepackage{enumerate}
\usepackage[left=3cm,right=2cm,top=2cm,bottom=2cm]{geometry}
\usepackage{epsfig}
 \renewcommand{\baselinestretch}{0.93}
 \thispagestyle{empty}
 \pagestyle{empty}
\setlength{\textheight}{53pc} \setlength{\textwidth} {36pc}
\setlength{\topmargin}{-4mm}
\usepackage{multicol}
\newcommand*{\email}[1]{%
    \normalsize\href{mailto:#1}{#1}\par
    }
\setlength{\columnsep}{1cm}
\newcommand{\N}{\mathbb{N}}
\newcommand{\calM}{\cal{M}}
\newcommand{\calP}{\cal{P}}
\newcommand{\F}{\mathbb{F}}
\newcommand{\R}{\mathbb{R}}
\newcommand{\Z}{\mathbb{Z}}
\newcommand{\Q}{\mathbb{Q}}
\newcommand{\C}{\mathbb{C}}
\newcommand{\bbH}{\mathbb{H}}


\theoremstyle{plain}
\newtheorem{thm}{Teorema}[section]
\newtheorem{prop}[thm]{Proposición}
\newtheorem{lemma}[thm]{Lema}
\newtheorem{cor}[thm]{Corolario}
\theoremstyle{definition}
\newtheorem{rec}[thm]{Recuerdo}
\newtheorem{defin}[thm]{Definición}
\newtheorem{ejem}[thm]{Ejemplo}
\newtheorem{ejer}[thm]{Ejercicio}
\newtheorem{sol}[thm]{Solución}
\newtheorem{rmk}[thm]{Observación}
\author{Marcos Morales}
\affil{Universidad Finis Terrae}
\title{Primera Semana\\}






\begin{document}


\begin{picture}(400, 75)
   \put(-40,15){\includegraphics[width=5cm]{UFT.png}}
\end{picture}


\begin{center}
\huge{D\'ecima Semana}
\end{center}

\begin{center}
\Large{ Marcos Morales, Camila Muñoz}
\end{center}

\begin{center}
	\large{Cálculo Vectorial}
\end{center}
\begin{center}
\setlength{\parindent}{0cm}\large{Clases centradas para capacitar a los estudiantes de ingeniería en Cálculo Vectorial. \\}
\end{center}


\vspace{1cm}

\begin{center}
	\large{Contenidos:}
\end{center}
\vspace{0.1cm}
\begin{center}
\begin{minipage}{0.5\textwidth}
\begin{enumerate}
    \item Integrales de L\'inea
    \item Campos conservativos
    \item  Teorema fundamental de las integrales de l\'inea
\end{enumerate}
\end{minipage}
\end{center}




\vspace{6cm}




\pagestyle{empty}


\setcounter{page}{1}
\pagestyle{plain}
\setlength{\parindent}{0cm}





\newpage

\section{Integrales de L\'inea}

Sea $C$ una curva param\'etrica, dada por $r(t)=(x_1(t),x_2(t),...,x_n(t))$, con $\alpha \leq t \leq \beta$. Existen dos tipos de definiciones para integrales de l\'inea.\\

\begin{defin}Si $f(x_1,x_2,...,x_n)$ es una funci\'on escalar, es decir, $f: \mathbb{R}^n \to \mathbb{R}$, entonces se define la integral de l\'inea de $f$ sobre $C$ como
\end{defin}

\begin{align*}
    \int_C f  ds = \int_\alpha^\beta f(x_1(t),x_2(t),...,x_n(t))||r'(t)|| dt
\end{align*}

\begin{defin}Si $f(x_1,x_2,...,x_n)$ es una funci\'on vectorial, es decir, $f: \mathbb{R}^n \to \mathbb{R}^m$, entonces se define la integral de l\'inea de $f$ sobre $C$ como
\end{defin}

\begin{align*}
    \int_C f \cdot dr = \int_\alpha^\beta f(x_1(t),x_2(t),...,x_n(t)) \cdot r'(t) dt
\end{align*}


Se puede demostrar en ambos casos, que el valor de la integral de l\'inea es independiente de la parametrizaci\'on de la curva $C$. \\

La raz\'on por la cual hay una diferencia importante en la definici\'on es porque las dos definiciones tienen interpretaciones muy diferentes. \\

La integral de l\'inea de un campo escalar positivo $f(x,y)$ representa el \'area lateral de un lado de la cortina cuya base es $C$ y cuya altura por encima de $(x,y)$ es $f(x,y)$

\begin{figure}[h]
    \centering
    \includegraphics[width=0.4\textwidth]{Imagenes/cortina.png}
    \caption{Interpretaci\'on geom\'etrica de la integral de l\'inea para campos escalares}
    \label{fig:etiqueta}
\end{figure}

Por otra parte, las integrales de l\'inea de los campos vectoriales tienen interpretaciones f\'isicas muy importantes, pues representan como se comporta una funci\'on a lo largo de una trayectoria, por ejemplo, el trabajo realizado por una fuerza $\Vec{F}$ a lo largo de una trayectoria $C$ viene dado por

\begin{align*}
    W = \int_C \vec{F} \cdot  d\vec{r}
\end{align*}

\newpage

\begin{ejem}Calcular la integral de línea del campo escalar
\end{ejem}
\[
f(x, y) = x^2 + y^2
\]
a lo largo de la curva $C$, que es el cuarto de circunferencia de radio $2$ en el primer cuadrante, desde $(2,0)$ hasta $(0,2)$. \\


La curva puede parametrizarse como

\[
r(t) = (2\cos t, 2\sin t), \quad 0 \leq t \leq \frac{\pi}{2}
\]

Ahora evaluamos en la funci\'on

\[
f(x(t), y(t)) = (2\cos t)^2 + (2\sin t)^2 = 4(\cos^2 t + \sin^2 t) = 4
\]

Por otra parte tenemos que

\begin{align*}
    \|r'(t)\| &= \|(-2\sin t, 2\cos t)\| \\
    &= \sqrt{4\sin^2 t + 4\cos^2 t} \\
    &= \sqrt{4} \\
    &= 2
\end{align*}

Por \'ltimo calculamos la integral



\begin{align*}
    \int_C f(x, y) \, ds &= \int_0^{\pi/2} f(r(t)) \cdot \|r'(t)\| \, dt \\
    &=  \int_0^{\pi/2} 4 \cdot 2 \, dt \\
    &= \int_0^{\pi/2} 8 \, dt \\
    &= 8 \cdot \frac{\pi}{2} \\
    &= 4\pi
\end{align*}






\begin{ejem}Sea el campo vectorial $F(x, y) = (y, x)$.
\end{ejem}
Calcular:

\[
\int_C F \cdot dr
\]
donde $C$ es el segmento de recta que une $(0, 0)$ con $(1, 1)$. \\


\begin{sol}Primero hay que parametrizar la curva $C$, tenemos que
\end{sol}

\[
r(t) = (t, t), \quad 0 \leq t \leq 1
\]
\[
r'(t) = (1, 1)
\]

Evaluamos la funci\'on en la curva

\[
\vec{F}(r(t)) = \vec{F}(t, t) = (t, t)
\]

Calculamos el producto punto

\begin{align*}
    \vec{F}(r(t)) \cdot r'(t) &= (t, t) \cdot (1, 1) \\
                         &= t + t \\
                         &= 2t \\
\end{align*}
Por \'ultimo calculamos la integral de l\'inea

\begin{align*}
    \int_C F \cdot dr &= \int_0^1 2t \, dt \\
&= \left[t^2\right]_0^1 \\
&= 1
\end{align*}



\begin{ejer}Sea el campo vectorial
\end{ejer}
\[
F(x, y) = (x^2 - y,\ xy + 2)
\]
Y sea \( C \) el semicírculo superior de radio 1, centrado en el origen, que va desde el punto \( (1, 0) \) hasta el punto \( (-1, 0) \). Calcular la integral de línea:

\[
\int_C F \cdot dr
\]

\begin{ejer}Sea el campo escalar
\end{ejer}
\[
f(x, y) = x + y
\]

Considere la recta \( C \)  que une los puntos \( (0, 0) \) y \( (1, 2) \). Calcular la integral de línea
\[
\int_C f(x, y) \, ds
\]

\newpage

\subsection{Conjuntos simplemente conexos}

La definici\'on formal de un conjunto simplemente conexo es bastante confusa \\

\begin{defin}Sea $X$ conjunto, $X$ se dice simplemente conexo si para toda aplicaci\'on continua $f:[0,1] \to X$ tal que $f(0)=f(1)=p$, con $p \in X$, existe una apliaci\'on continua $H: [0,1] \times [0,1] \to X$ tal que $H(s,0) =f(s)$ y $H(s,1)=p$. \\
\end{defin}

La idea intuitiva de esta definici\'on es que se est\'a deformando la funci\'on $f$ de forma continua hasta llegar a un punto.\\

Informalmente, un objeto es simplemente conexo si está formado por una sola pieza y no contiene agujeros que lo atraviesen.


\begin{figure}[h]
    \centering
    \includegraphics[width=0.4\textwidth]{Imagenes/conexo.png}
    \caption{Este espacio no es simplemente conexo. Aunque es conexo por caminos, tiene agujeros que lo atraviesan.}
    \label{fig:etiqueta}
\end{figure}


En general no vamos a usar la definici\'on ni tampoco vamos a pedir demostrar cuando un conjunto es simplemente conexo, pero usaremos el concepto en teoremas y resultados.






\newpage

\subsection{Campos conservativos}

\begin{defin}Sea $F:A \subseteq \mathbb{R}^n \to \mathbb{R}^n$, donde $A$ es un conjunto abierto simplemente conexo , se dice que $F$ es \textbf{conservativo} si existe un campo escalar $f:\mathbb{R}^n \to \mathbb{R}$ tal que
\end{defin}

\begin{align*}
    \nabla f = F
\end{align*}

A la funci\'on $f$ tal que $\nabla f = F$, se le llama \textbf{funci\'on potencia}. \\

Existen formas muy eficientes para saber si un campo vectorial es conservativo en dos o tres dimensiones. \\

\begin{prop}Supongamos que $F: A \subseteq \mathbb{R}^2 \to \mathbb{R}^2$, con $A$ simplemente  conexo, $F(x,y)=M(x,y)\hat{\imath}+N(x,y)\hat{\jmath}$, entonces $F$ es conservativo  si y solo si
\end{prop}

\begin{align*}
    M_y = N_x
\end{align*}

Para tres dimensiones necesitamos el concepto de rotacional. \\


\begin{defin}Sea $F: A \subseteq \mathbb{R}^3 \to \mathbb{R}^3$ , $F(x,y,z)=M(x,y,z)\hat{\imath}+N(x,y,z)\hat{\jmath}+C(x,y,z)\hat{k}$, se define el \textbf{rotacional} de $F$ como
\end{defin}

\begin{align*}
    \text{rot} ( F ) &= \begin{vmatrix}
        \hat{\imath} & \hat{\jmath} & \hat{k} \\
        \frac{\partial}{\partial x} & \frac{\partial}{\partial y} & \frac{\partial}{\partial z} \\
        M & N & C \\
    \end{vmatrix} \\
    &= (C_y-N_z)\hat{\imath}+ (M_z-C_x)\hat{\jmath}+ (N_x-M_y)\hat{k}
\end{align*}

\begin{prop}Supongamos que $F: A \subseteq \mathbb{R}^3 \to \mathbb{R}^3$ con $A$ simplemente conexo, $F(x,y,z)=M(x,y,z)\hat{\imath}+N(x,y,z)\hat{\jmath}+C(x,y,z)\hat{k}$, entonces $F$ es conservativo  si y solo si
\end{prop}

\begin{align*}
    \text{rot} ( F ) = (0,0,0)
\end{align*}


\begin{ejem}Determinar si el campo $F(x,y) = (6xy^3,9x^2y^2)$ es conservativo o no \\
\end{ejem}

\begin{sol}Tenemos que $M= 6xy^3$ y $ N = 9x^2y^2$, luego
\end{sol}

\begin{align*}
    M_y &= 18xy^2 \\
    N_x &= 18xy^2 \\
\end{align*}

Por lo tanto el campo es conservativo. \\

\begin{ejem}Determinar si el campo $F(x,y,z) = (2xyz^3,x^2z^3, 3x^2yz^2)$ es conservativo o no \\
\end{ejem}

\begin{sol}Tenemos que $M= 2xyz^3$, $ N = x^2z^3$ y $C=3x^2yz^2 $, luego
\end{sol}

\begin{align*}
     \text{rot} ( F )  &= (C_y-N_z)\hat{\imath}+ (M_z-C_x)\hat{\jmath}+ (N_x-M_y)\hat{k} \\
     &= (3x^2z^2-3x^2z^2)\hat{\imath}+(6xyz^2-6xyz^3)\hat{\jmath}+(2xz^3-2xz^3)\hat{k} \\
     &= (0,0,0)
\end{align*}

Por lo tanto el campo es conservativo


\newpage

Ahora, si bien podemos determinar cuando un campo es o no conservativo falta poder encontrar la funci\'on $f$ tal que $\nabla f = F$, para ello vamos a usar el siguiente teorema cuya demostraci\'on omitiremos \\

\begin{thm}Sea $F(x,y,z)$ campo conservativo, considere $C$ la curva parametrizada por $\lambda(t)=(tx,ty,tz)$ con $t \in [0,1]$, consideramos
\end{thm}

\begin{align*}
    f(x,y,z) = \int_C F \cdot d\lambda
\end{align*}

Entonces $\nabla f = F$, en otras palabras $f$ es funci\'on potencia de $F$. \\

Este teorema funciona tanto para funciones en dos variables como en tres variables. \\

\begin{ejem}Sea $F(x,y) =(x+y^2,2xy)$, demuestre que $F$ es conservativo y luego calcule su funci\'on potencia \\
\end{ejem}

\begin{sol}Notamos que $M_y=N_x= 2y$ por lo tanto el campo es conservativo, la funci\'on potencia ser\'a
\end{sol}

\begin{align*}
    \int_C F \cdot d\lambda &= \int_0^1 F(tx,ty) \cdot \lambda'(t)dt \\
    &= \int_0^1 (tx+t^2y^2,2t^2xy) \cdot (x,y)dt \\
    &= \int_0^1 tx^2+t^2xy^2 + 2t^2xy^2dt \\
    &= \int_0^1 tx^2+3t^2xy^2dt \\
    &= \frac{x^2t^2}{2}+xy^2t^3 \bigg|_0^1\\
    &= \frac{x^2}{2}+xy^2
\end{align*}

Por lo tanto la funci\'on potencia es $f(x,y) = \displaystyle \frac{x^2}{2}+xy^2$


\begin{ejer}Sea $F(x,y,z) =(3y^2z+ye^x,6xyz+e^x,3xy^2)$, demuestre que $F$ es conservativo y luego calcule su funci\'on potencia \\
\end{ejer}

\begin{sol}Tenemos que
\end{sol}

\begin{align*}
     \text{rot} ( F )  &= (C_y-N_z)\hat{\imath}+ (M_z-C_x)\hat{\jmath}+ (N_x-M_y)\hat{k} \\
     &= (6xy-6xy)\hat{\imath}+(3y^2-3y^2)\hat{\jmath}+(6yz+e^x-6yz-e^x)\hat{k} \\
     &= (0,0,0)
\end{align*}

Por lo tanto $F$, es conservativo, la funci\'on potencia ser\'a

\begin{align*}
    \int_C F \cdot d\lambda &= \int_0^1 F(tx,ty,tz) \cdot \lambda'(t)dt \\
    &= \int_0^1 (3t^3y^2z+tye^{tx},6t^3xyz+e^{tx},3t^3xy^2) \cdot (x,y,z)dt \\
    &= \int_0^1 3t^3y^2zx+txye^{tx} + 6t^3xy^2z+ye^{tx}+ 3t^3xy^2zdt \\
    &= \int_0^1 12t^3y^2zx+ye^{tx}  + txye^{tx} dt\\
    &= 3t^4y^2zx+\frac{y}{x}e^{tx}+yte^{tx}-\frac{ye^{tx}}{x} \bigg|_0^1 \\
    &= 3y^2zx+\frac{y}{x}e^x+ye^{x}-\frac{ye^{x}}{x} \\
    &= 3xy^2z+ye^x
\end{align*}

Por lo tanto la funci\'on potencia es $f(x,y,z) = 3xy^2z+ye^x$


\newpage


\subsection{Teorema fundamental de las integrales de l\'inea}

El teorema fundamental de las integrales de linea es an\'alogo al teorema fundamental del c\'alculo\\

\textbf{ Teorema fundamental de las integrales de linea: }Sea $C$ una curva suave parametrizada por $r(t)$, con $a \leq t \leq b$, contenida en una región abierta $R$ simplemente conexa.   Si  $F = M\hat{\imath} + N\hat{\jmath}$  es conservativo en $R$, con funci\'on potencia $f$,  y sus campos escalares $M, N$ son continuas en $R$, entonces

\begin{align*}
    \int_C F \cdot dr =  \int_C \nabla f \cdot dr =f(r(b))-f(r(a))
\end{align*}

\begin{ejem}Sea $F(x,y)= (2xy+3,x^2+\cos(y))$, usando el teorema fundamental de las integrales de línea, calcule
\end{ejem}

\begin{align*}
    \int_C F \cdot dr
\end{align*}

Donde $C$ es cualquier curva suave que va desde el punto $A=(0,0)$, hasta el punto $B=(2,1)$. \\

\begin{sol}No parematrizamos la curva, primero vemos que sea conservativo, tenemos que $M= 2xy+3$ y $N=x^2+\cos(y)$, luego $M_y=N_x=2x$, por lo tanto es conservativo, la función potencia será
\end{sol}

\begin{align*}
    f(x,y) &= \int_0^1 F(tx,ty) \cdot (x,y) dt \\
    &= \int_0^1 (2xyt^2+3,x^2t^2+\cos(ty)) \cdot (x,y) dt \\
    &= \int_0^1 2x^2yt^2+3x +x^2t^2y+y\cos(ty)) dt \\
    &= \int_0^1 3x^2yt^2+3x +y\cos(ty)) dt \\
    &= \left( 3x^2y \frac{t^3}{3}+3xt +\sin(ty)\right)\Big|_0^1 \\
    &= x^2y+3x+\sin(y)
\end{align*}

Luego por el teorema fundamental de las integrales de línea el valor de la integral es independiente de la parametrización de la curva que se recorre, por lo tanto


\begin{align*}
    \int_C F \cdot dr &= f(2,1)- f(0,0) \\
    &= 10+\sin(1) \\
\end{align*}



\end{document}